\documentclass[12pt]{article}
\usepackage{amsmath, amssymb}
\usepackage{geometry}
\geometry{margin=1in}

\title{Lecture Scribe\\
CSE400 -- Fundamentals of Probability in Computing\\
Lecture 10: Randomized Min-Cut Algorithm}
\author{Instructor: Dhaval Patel, PhD}
\date{February 5, 2026}

\begin{document}
\maketitle

\section{Introduction / Scope}

This lecture covers the \textbf{Min-Cut Problem}, including:

\begin{itemize}
\item Motivation and applications
\item Definition of cut-set and minimum cut
\item Successful and unsuccessful min-cut runs
\item Max-Flow Min-Cut Theorem
\item Deterministic Min-Cut Algorithm (Stoer--Wagner Algorithm)
\item Randomized Min-Cut Algorithm (Karger’s Algorithm)
\item Comparison between deterministic and randomized approaches
\item Theorem on probability of success of randomized min-cut
\end{itemize}

This scribe strictly follows the lecture slides. No intuition, commentary, simplifications, or additional material beyond the lecture content is included.

\section{Min-Cut Problem}

\subsection{Why Use Min-Cut?}

We use the min-cut algorithm in various applications to solve problems related to network connectivity, reliability, and optimization.

\subsubsection*{Applications}

\begin{enumerate}
\item \textbf{Network Design} \\
Min cut helps in improving the efficiency of communication and optimizing network flow. The algorithm is used in network design to find the minimum capacity cut.

\item \textbf{Communication Networks} \\
For understanding the vulnerability of the networks to failures minimum cut can be useful. It helps building robust and fault-tolerant communication networks.

\item \textbf{VLSI Design} \\
In Very Large Scale Integration (VLSI) design, the algorithm is useful for partitioning circuits into smaller components leading to reduced interconnectivity complexity.
\end{enumerate}

Reference mentioned in lecture: \\
Section 1.5, \textit{Application: A Randomized Min-Cut Algorithm} in \textit{Probability and Computing (2nd Edition)}.

\subsection{What is Min-Cut?}

\subsubsection*{Definition: Cut-Set}

A cut-set in a graph is a set of edges whose removal breaks the graph into two or more connected components.

\subsubsection*{Definition: Minimum Cut}

Given a graph
\[
G = (V, E)
\]
with \( n \) vertices, the minimum cut --- or min-cut --- problem is to find a minimum cardinality cut-set in \( G \).

\subsubsection*{Remark (From Lecture)}

Min-cut algorithms, such as Karger’s algorithm, are random and can be sensitive to the initial choice of edges. If the algorithm happens to contract critical edges early in the process, it might find a smaller cut.

\subsection{Edge Contraction}

The main operation in the algorithm is \textbf{edge contraction}, which removes an edge from a graph while simultaneously merging the two vertices.

\subsubsection*{Edge Contraction Operation}

In contracting an edge \( (u, v) \):

\begin{itemize}
\item Merge the two vertices \( u \) and \( v \) into one vertex.
\item Eliminate all edges connecting \( u \) and \( v \).
\item Retain all other edges in the graph.
\item The new graph may have parallel edges but no self-loops.
\end{itemize}

\section{Successful and Unsuccessful Min-Cut Runs}

\subsection{Successful Min-Cut Run}

A successful min cut run refers to the success in the outcome of an algorithm designed to find the minimum cut in a graph.

(Figure: Successful Run — as shown in lecture slide)

\subsection{Unsuccessful Min-Cut Run}

An unsuccessful min-cut run refers to an iteration of a min-cut algorithm where the algorithm fails to correctly identify the minimum cut of a given graph.

(Figure: Unsuccessful Run — as shown in lecture slide)

\section{Max-Flow Min-Cut Theorem}

\subsection*{Statement of Theorem}

\begin{quote}
In a flow network, the maximum amount of flow passing from the source to the sink is equal to the total weight of the edges in a minimum cut.
\end{quote}

\subsection*{Definitions}

\textbf{Capacity of a Cut}

Capacity of a cut = the sum of the capacity of the edges in the cut that are oriented from a vertex \( \in X \) to a vertex \( \in Y \).

\medskip

\textbf{Minimum Cut}

Minimum cut = the cut in the network that has the smallest possible capacity.

Minimum cut capacity = the capacity of the minimum cut.

\medskip

\textbf{Maximum Flow}

Maximum flow = the largest possible flow from source \( S \) to sink \( T \).

\section{Deterministic Min-Cut Algorithm}

\subsection{Stoer--Wagner Min-Cut Algorithm}

Let \( s \) and \( t \) be two vertices of a graph \( G \). \\
Let \( G/\{s, t\} \) be the graph obtained by merging \( s \) and \( t \).

Then a minimum cut of \( G \) can be obtained by taking the smaller of:

\begin{itemize}
\item A minimum \( s\text{-}t \)-cut of \( G \), and
\item A minimum cut of \( G/\{s, t\} \).
\end{itemize}

\subsubsection*{Justification (As Stated in Lecture)}

The theorem holds since:

\begin{itemize}
\item Either there is a minimum cut of \( G \) that separates \( s \) and \( t \), then a minimum \( s\text{-}t \)-cut of \( G \) is a minimum cut of \( G \);
\item Or there is none, then a minimum cut of \( G/\{s, t\} \) does the job.
\end{itemize}

\subsection{Pseudocode}

\subsubsection*{Algorithm 1: MinimumCutPhase(G, a)}

\begin{enumerate}
\item \( A \leftarrow \{a\} \)
\item while \( A \neq V \) do
\begin{enumerate}
\item add to \( A \) the most tightly connected vertex;
\end{enumerate}
\item return the cut weight as the ``cut of the phase'';
\end{enumerate}

\subsubsection*{Algorithm 2: MinimumCut(G)}

\begin{enumerate}
\item while \( |V| \ge 1 \) do
\begin{enumerate}
\item choose any \( a \) from \( V \);
\item MinimumCutPhase(G, a);
\item if the cut-of-the-phase is lighter than the current minimum cut then
\begin{enumerate}
\item store the cut-of-the-phase as the current minimum cut;
\end{enumerate}
\item shrink \( G \) by merging the two vertices added last;
\end{enumerate}
\item return the minimum cut;
\end{enumerate}

\section{Randomized Min-Cut Algorithm}

\subsection{Why Randomized Algorithm?}

Randomized algorithms provide a probabilistic guarantee of success. It provides a more accurate estimate of the minimum cut with fewer iterations.

Additional points listed in lecture:

\begin{itemize}
\item Efficiency
\item Parallelization
\item Approximation Guarantees
\item Avoidance of Worst-Case Instances
\item Heuristic Nature
\item Robustness
\end{itemize}

\subsection{Karger’s Randomized Algorithm}

(Title slide introducing Karger’s algorithm; details correspond to randomized contraction method described earlier.)

\subsection{Comparison: Deterministic vs Randomized Min-Cut}

Any specific problem is responsible for deciding the approach.

\subsubsection*{Deterministic Min Cut}

\begin{itemize}
\item It always guarantees an exact minimum cut.
\item It may have higher time complexity for large graphs.
\item The Stoer-Wagner algorithm has time complexity:
\end{itemize}

\[
O(V \cdot E + V^2 \log V)
\]

where \( V \) is the number of vertices and \( E \) is the number of edges.

\subsubsection*{Randomized Min Cut}

\begin{itemize}
\item It provides an approximate minimum cut with high probability.
\item Karger’s algorithm has time complexity:
\end{itemize}

\[
O(V^2)
\]

where \( V \) is the number of vertices in the graph.

\section{Theorem for Min-Cut Set}

\[
\text{The algorithm outputs a min-cut set with probability at least } \frac{2}{n(n-1)}.
\]

\section{Python Simulation}

Instruction from lecture:

\begin{itemize}
\item Open Campuswire post regarding today’s lecture (L10).
\item Download the \texttt{.ipynb} file pasted in the comments.
\end{itemize}

\bigskip
\textbf{End of Lecture 10}

\end{document}
